% P8 introduction: the deployment question
\section{Introduction}
\label{sec:p8-intro}

A bank that already runs a tuned gradient-boosted tree over its card-transaction
stream faces a concrete deployment question when the promotional material arrives
promising ``AI-powered fraud detection'': what, exactly, would a large language
model (LLM) add? The incumbent is formidable. Gradient-boosted decision trees
remain the strongest general-purpose learners on medium-sized tabular data
\citep{grinsztajn2022tree, shwartzziv2021tabular}, and on our own custom fraud
benchmark a stock XGBoost \citep{chen2016xgboost} configuration recorded a
held-out AUC of $0.7955$ in the current quick artifact-backed run
(\texttt{experiments/results/quick\_20260704/qp8\_fraud.tsv}). Any
proposal to replace it must beat not just that number but the operational
envelope around it: single-digit-millisecond scoring, near-zero marginal cost
per transaction, monotone and calibrated probability outputs that downstream
rules consume directly, and an input surface that an adversary cannot talk to.

This paper takes the question seriously rather than rhetorically. We fine-tune
an LLM on textual serializations of the same training table---the recipe shown
by TabLLM to be surprisingly competitive in few-shot regimes
\citep{hegselmann2023tabllm}---and it scores below the tree in that
head-to-head: the current Qwen3.5-4B SFT run reaches accuracy $0.792$ but AUC
$0.48268$ on a 500-row positive-enriched held-out evaluation
(\secref{sec:p8-scorer}). We then argue
that this gap is the least interesting fact in the study. Scoring
fixed-schema transactions is the one task in the fraud pipeline that is
\emph{already solved}; the interesting question is which tasks in that pipeline
were previously unaddressable by any model, and whether the LLM addresses them.

We find four such tasks, and none of them is tabular scoring
(\secref{sec:p8-taxonomy}). A vision--language model can read the pixels and
text of a disputed receipt or a doctored identity document---inputs that no
feature pipeline delivers to a tree \citep{xu2024fakeshield,
wang2025forensicsbench}. An LLM can draft the Suspicious Activity Report (SAR)
narrative that FinCEN requires and the adverse-action reasons that
ECOA/Regulation~B requires \citep{fincen2003sarnarrative, cfr1002_9}---a
regulated text-generation category in which a tree cannot compete because it
does not produce text at all. An LLM can triage a genuinely novel fraud
typology from a description and a handful of examples, before any labels exist
for supervised training. And an LLM agent can work the alert queue---gathering
context, checking consistency, drafting dispositions---at roughly $85\times$
lower cost than a human analyst by our internal estimate.

Our constructive contribution is the resulting division of labor
(\secref{sec:p8-architecture}): the LLM as \emph{sensor} (converting
unstructured evidence into features the tree can score) and \emph{scribe}
(converting scores and evidence into regulated narratives), with XGBoost
retained as the \emph{scorer} and an agentic triage loop downstream of the
score. We ground the scribe role in the primary regulatory texts
(\secref{sec:p8-compliance}), because the compliance argument is the one that
survives even if future LLMs close the tabular accuracy gap.

\paragraph{Program context.} This study is the deliberate side-probe of a
reproducibility program whose thesis is ``measure capability, don't trust the
label.'' The main program applies that discipline to reinforcement-learning
post-training claims; here we apply it to the label ``AI'' itself. ``AI'' in a
fraud-detection pitch bundles at least five distinct capabilities---tabular
scoring, perception, narration, few-shot generalization, and agency---and they
do not rise or fall together. Measured separately, one of them (scoring)
favors a 2016-era tree in the current quick artifacts, and we argue the other four
favor the LLM. Conflating them produces both failure
modes we see in practice: buying an LLM to do a tree's job, and refusing an LLM
for jobs no tree can do.

This side study is intentionally parked outside the ZVF thesis. It contributes
an example of the program's measurement discipline, but supplies no evidence for
GRPO signal starvation, group-size control, PPO, or SAO. Its scientific claims
must stand on the fraud artifacts and cited domain literature alone.

\paragraph{Contributions.}
(i) An honest head-to-head between XGBoost and a fine-tuned LLM on a custom
fraud dataset, with the tree keeping the scorer seat on our internal accuracy
records, the current quick artifacts, \emph{and} on
latency, cost, calibration, and prompt-injection exposure;
(ii) a four-way capability-gap taxonomy of where the LLM genuinely adds value,
each gap being a task category rather than an accuracy delta;
(iii) a hybrid sensor/scorer/scribe architecture with post-score agentic
triage; and
(iv) a compliance grounding of the scribe role in FinCEN SAR requirements and
ECOA/Regulation~B adverse-action requirements, citing primary sources.
The older $0.975$/$0.948$ pair survives only as a historical internal record;
the scorer claims below use the current quick artifacts unless explicitly
marked otherwise. \secref{sec:p8-limitations} states the boundaries
explicitly.
